\documentclass[acmsmall,screen]{acmart}

%% ---- FSE 2026 / PACMSE (Proceedings of the ACM on Software Engineering) ----
%% Named (non-anonymous) version of the focused paper (source: fse_focused_v2.md).
%% v3: adds Section 9, a separately and later pre-registered live study that tests
%% the detection-cost coupling on a second model family (GPT-5.5). No change to the
%% original V1->V2 narrative or its frozen numbers; see the Integrity log (D-V3-1).
\settopmatter{printacmref=false, printccs=false}
\setcopyright{none}
\acmJournal{PACMSE}
\acmVolume{1}
\acmNumber{FSE}
\acmArticle{1}
\acmYear{2026}
\acmMonth{7}

\usepackage{booktabs}
\usepackage{tabularx}
\usepackage{enumitem}
\usepackage{listings}

\lstset{basicstyle=\ttfamily\footnotesize, breaklines=true,
  breakatwhitespace=false, columns=fullflexible, keepspaces=true,
  xleftmargin=4pt, aboveskip=6pt, belowskip=6pt, frame=none,
  extendedchars=true, inputencoding=utf8, literate={§}{{\S}}1}

\begin{document}

\title{When the Monitor Blinds Itself: Failure Detection in Multi-Agent LLM
Systems Is Limited by Observation---and Coupled to Its Cost}

\author{Anshuman Mullick}
\affiliation{%
  \institution{Bilkent University}
  \city{Ankara}
  \country{Turkey}
}
\email{anshumanmullick7@gmail.com}

\author{Eray T\"uz\"un}
\affiliation{%
  \institution{Bilkent University}
  \city{Ankara}
  \country{Turkey}
}
\email{eraytuzun@cs.bilkent.edu.tr}

\begin{abstract}
A multi-agent LLM system runs many workers in parallel against a single plan.
When that plan quietly breaks mid-run, the workers keep spending tokens on a dead
plan until the very end. We built a monitor to catch it, wrote down the bar it
had to clear in advance, and ran it to a verdict we were not free to renegotiate.
It failed---three times---and the failure is the finding.

The finding has two halves. First, catching a broken plan is limited by what the
monitor gets to \emph{see}, not by how clever it is---and a monitor that
interrupts too eagerly spends on false alarms the very budget a second look
needs, so accuracy and restraint draw on one finite pot. Second, the only
mechanism that restores that second look---re-reading the plan and writing fresh
checks---is also the system's largest cost, so cutting the cost risks the
detection it buys: \textbf{detection and cost are coupled.}

Our first version failed badly: 35\% detection, a ``judge'' stage that
rubber-stamped its own false alarms, and a cost worse than the plain baseline it
was meant to save---which a \emph{naive} baseline then beat outright, at zero
false alarms. A trace-level diagnosis, reviewed by two outside red teams, found it
had starved its own observation. The redesign restores observation directly and
catches 24 of 31 injected faults at zero false alarms---but at a 55.5\% clean
overhead against a 12\% cap, an \textbf{overall fail} on cost, because that
overhead \emph{is} the plan-reading step that does the detecting, and spreading it
over more workers only widens the gap. A separately and later pre-registered study
then put a second vendor's model (GPT-5.5) in the check-writer's seat---a cheaper,
better writer---and the bar failed a third time: shortening the checklist to save
money drops exactly the faults that need a specific probe and \emph{destabilizes}
detection. The coupling holds across both model families.
\end{abstract}

\keywords{multi-agent LLM systems, runtime verification, self-adaptive
systems, failure detection, interrupt policy, pre-registration, agent
orchestration}

\maketitle

\section{Introduction}
Twelve seconds before the fault it existed to catch, our monitoring system
destroyed its own coverage.

Here is what happened in one run. A check fired on a perfectly healthy file
listing. A ``judge'' model approved the alarm as real, at 0.98 confidence, by
hallucinating a misreading of the listing. The orchestrator paused its workers
and replanned. The new plan got a fresh set of checks---and those checks
covered none of the surfaces the old ones had watched. Twelve seconds later the
real fault fired into a world no one was watching. The run finished, confident
and wrong.

Every stage did its job. The system as a whole blinded itself.

This paper is about that mechanism. In a multi-agent system on a fixed budget,
catching a broken plan is bounded by \emph{observation}, not by the monitor's
\emph{intelligence}---and a monitor that is too eager starves its own
observation. We establish this through three acts: a pre-registered experiment
that killed most of our design, a careful diagnosis of why, and a redesign that
puts observation back at the center.

\textbf{The stakes.} Multi-agent LLM systems are expensive. They use, on the
order of 15$\times$ the tokens of single-agent chat, as one production report
puts it \cite{ref27}. When a plan
breaks mid-run, plain batch orchestration only finds out at the end---after
every parallel worker has spent that multiplier on a dead plan. These failures
are real and well-catalogued \cite{ref14, ref16}. And the \emph{mechanism} to
interrupt a run already exists in production frameworks \cite{ref22}. What is
missing is the \emph{policy}: something that decides what to watch for, whether
an anomaly actually breaks the plan, and when interrupting is worth it.

\textbf{What we built.} We proposed a layer called Sentinel Protocol. The idea
is simple. Before the workers run, an LLM reads the plan once and writes out a
set of small, concrete checks---we call them \emph{tripwires}---that say what
each step is trusting about the world. The workers then watch for those checks
cheaply, by pattern matching, and a confirmed break interrupts the orchestrator
so it can replan. We pay for expensive reasoning once, up front, and keep the
running system cheap. (For software-engineering readers: it is a classic
self-adaptive control loop \cite{ref23}, but its monitor is \emph{generated
from the plan at runtime} rather than written by hand.)

\textbf{What we did to it.} We then did something the field rarely does to its
own proposals. We wrote down pass/fail lines \emph{in advance}---with the
consequence of each failure decided ahead of time---froze them, and ran the
experiment to a verdict we could not move after seeing the data.

The verdict killed most of the design. Detection was 35\%, against a 60\% bar
and a 40\% floor below which the claim died outright. The ``judge'' stage
approved all 34 alarms it ever saw, including all 18 false ones---a rubber
stamp, not a filter. The cost was worse than the plain baseline it was built to
save. One claim survived: where a fault keeps showing up on a surface someone
watches, our checks catch it about 3$\times$ faster than a cost-matched
periodic re-check (median 3 vs.\ 9 tool calls).

\textbf{Why it failed.} A failed verdict is only useful if you find out why. We
replayed the fully recorded runs, sent the diagnosis to two outside red teams
for blind review, and settled every testable claim with a deterministic replay
that cost no LLM calls. The diagnosis: detection here is observation-bounded,
and in our case the bound was self-inflicted. On identical seeds, ordinary
baseline systems re-looked at every starved surface (12 of 12)---but ours did
not. Its own false alarms had eaten the part of the run in which a second look
would have happened. Under a fixed budget, accuracy and restraint are not
separate dials.

\textbf{The redesign.} The fix follows from the diagnosis: stop relying on
workers passively noticing things, and instead have the monitor \emph{re-read
the world itself} with cheap, deterministic checks. We validated this mechanism
in offline replay before building it. Then we re-measured the new version
against a fresh set of pre-registered bars, frozen before any new data existed,
with two fault types whose exact settings were sealed away with someone outside
the project---so the design could not have seen them.

\textbf{The new verdict.} The redesign passed detection and failed cost. It
caught 24 of 31 injected faults, at \emph{zero} false alarms on clean runs,
recovering exactly the quiet, hard-to-see fault types the first version was
blind to. It even transferred partway to the sealed fault types it had never
seen: it caught the held-out RESOURCE\_BUDGET type 3 of 5 times, where every
passive baseline caught it 0 of 5---a \emph{directional} edge at $n=5$, not yet a
significant one (\S7). (The other held-out type, DEPENDENCY\_VERSION, was missed
by \emph{every} system, ours included---the same observation bound met on the
read side: a fault that fires before the monitor's first look cannot be caught by
comparison, by any system (\S7).) But its extra cost on clean runs was 55.5\%,
against a 12\% cap---an \textbf{overall fail}, driven entirely by cost.

So the thesis sharpens. Restoring observation fixes detection and ends the
self-harm. But the restored observation also has to be \emph{cheap}, and ours
is not.

\textbf{This is not a system that failed at its only job.} Where signals
recur, it detects about 3$\times$ faster than the obvious alternative. It
injures \emph{zero} clean runs in the noiseless mock (Section~10 measures the deployment
bound under benign noise), where both its own predecessor and the naive
baseline raise false alarms. And it carries part of its skill to fault types it
was never shown. The cost failure is not the absence of value. It is that the
value does not yet clear the cost of looking. Section 8 shows that this cost
failure and the detection success are two views of the same fact: you cannot
monitor a plan without reading it, and reading the plan is the irreducible cost.
Section 9 shows the same fact with a different vendor's model in the
check-writer's seat, under a separate frozen bar committed later---a cheaper,
better writer, and the coupling holds anyway, now because \emph{coverage} of the
plan is both what costs and what catches.

\textbf{Contributions.}
\begin{itemize}
  \item \textbf{The finding} (\S6, \S8, \S9): in budget-bounded multi-agent
  execution, detection is limited by observation, that limit can be
  self-inflicted, and---because the one thing that restores observation, reading
  the plan to write checks, is also the main cost---detection and cost are
  \emph{coupled}. The negatives are not a string of missed targets; they converge
  on one fact. That fact reproduces with a second vendor's model in the
  check-writer's seat (\S9), through a different mechanism---checklist coverage,
  not reasoning effort: the coupling is not one vendor's artifact.
  \item \textbf{A pre-registration method for agent-systems research} (\S5,
  \S7, \S9) that actually bound us: frozen bars returned \emph{fail} \emph{three}
  times---twice on both versions of our own system, and a third time in a
  separately and later pre-registered live study of a different vendor's model
  (\S9)---and the pre-decided consequences ran as written each time. We add
  binding outside red-team review, a data embargo, a log of every run, sealed
  held-out fault types, and a check-writer that is never told the fault list.
  \item \textbf{A diagnosis method} (\S6): deterministic replay over fully
  recorded runs, used to settle outside criticism with evidence rather than
  argument, and to test the redesign before building it.
  \item \textbf{TripwireBench} (\S4): a benchmark of injected faults that are
  \emph{verified to actually break a plain baseline}, with clear recovery
  labels, sealed held-out types, and cost/speed/noise metrics.
  \item \textbf{The surviving result and the naive baseline} (\S5, \S7): a
  3$\times$ speed edge where signals recur, and a no-frills baseline that
  caught 56\% in the pilot at zero false alarms and beat our full first version in every
  category---the honest floor any successor has to clear.
\end{itemize}

\textbf{Three questions organize the paper.} \emph{RQ1 (the observation bound):}
what does a plan-compiled tripwire monitor detect that a passive baseline cannot,
and what limits its recall? (\S5, \S7) \emph{RQ2 (the cost coupling):} does the
monitor's detection value clear its cost, or are the two coupled? (\S8)
\emph{RQ3 (generality):} does that coupling hold across model families, and does
the noiseless false-alarm floor survive benign deployment noise? (\S9, \S10)

\section{Background and Related Work}
\textbf{Self-adaptive systems.} Software-engineering readers will recognize the
shape: a monitor--analyze--plan--execute control loop, a staple of
self-adaptive systems for two decades \cite{ref23, ref24}, including recent work
that puts LLMs inside the loop \cite{ref25, ref26}. We do not claim the loop.
What is new is where its parts come from: the monitoring logic is \emph{written
from the plan, at runtime, by an LLM} rather than designed by hand. Our
experiment then shows that the analysis step is the hard part---an uncalibrated
filter does not just fail to help; it does damage (\S5).

\textbf{The nearest prior work, confronted.} The closest ancestor is RVPLAN
\cite{ref6, ref34}, which turns a planner's stated preconditions into runtime
monitors and replans when one breaks. ``Watch the plan's assumptions, interrupt
to replan'' is therefore established. But RVPLAN lives in a setting where three
things hold that do not hold for us:

\begin{itemize}
  \item \emph{The specification already exists.} RVPLAN translates
  preconditions that sit, machine-readable, in a planning file. Our checks have
  to be \emph{invented} from a plan written in natural language. Coverage is
  therefore bounded by what the writer can infer (\S6).
  \item \emph{The world narrates itself.} RVPLAN reads a clean stream of
  labeled facts in its own vocabulary. We read raw, unlabeled tool output,
  often only once---and the gap between ``watching'' and ``actually seeing'' is
  our central result.
  \item \emph{The stream is clean.} In a labeled stream a precondition simply
  holds or it does not; nothing can false-fire, so nothing needs filtering. In
  our setting the filtering layer is the whole game, and it is where our first
  version broke.
\end{itemize}

An established paradigm meets the LLM-agent setting and breaks in three
measurable ways. This paper is about what it takes to survive there.

\textbf{Other lines.} Runtime verification and monitor synthesis \cite{ref4,
ref5} give us the spec-to-monitor tradition; we trade full temporal logic for a
small predicate language over tool calls, and our complaint is about
\emph{passivity}, not the trade (\S3). LLM-agent guardrails \cite{ref9, ref10,
ref11, ref12, ref13} enforce externally written \emph{safety} rules; our checks
come from the plan, target plan validity and cost, and feed replanning. Failure
taxonomies \cite{ref14, ref20} catalog things agents do wrong; we cover the
complement---things the \emph{world} does that the plan didn't expect---and we
use them to write checks before the run, not to label failures after.
Observability tools \cite{ref15, ref16, ref17, ref18} explain failures after
the fact, for humans; we act on the same signals in flight. LangGraph's
interrupt \cite{ref22} is mechanism without policy; schedulers \cite{ref19}
consume interrupts once someone decides to raise them. We sit in the gap:
deciding \emph{what to watch} and \emph{whether to pull the trigger}.

\section{The System Under Test}
This section describes the system, but the system is our \emph{instrument}, not
our contribution. It is the thing we pre-registered, killed, diagnosed, and
rebuilt in order to reach the finding. We describe the current (v2) version and
note where the first version (v1) differed. Every change from v1 to v2 cites
the trace evidence that forced it (\S6), and v1's verdict is reported in full
no matter how v2 turns out (\S5, \S7).

\textbf{Three roles.} The \emph{orchestrator} owns the plan and is the only
part allowed to replan; it acts only on confirmed breaks. The \emph{sentinel}
runs once, up front: it reads the plan and writes the checks. The \emph{workers}
carry those checks and watch for them cheaply while they work, raising a
structured alarm on a match. The cycle is: plan, write checks, run, raise an
alarm, confirm-and-interrupt, replan. Pay for the thinking once; keep the
running system cheap and deterministic. The pilot added one rule the new
version follows: the running system also has to \emph{look}---not just wait to
be told.

\textbf{Checks and probes.} A check is concrete. It names a specific thing to
watch (a status code, a field, a structure, a hash), a specific value, and what
to do if it trips. Vague checks (``if the API changed'') are not allowed. One
rule we learned the hard way and now enforce: every constraint the check-writer
must follow has to live in a place the model actually reads---a field
description in the schema---because rules written in comments or prose simply do
not reach it.

The pilot showed the checks were fine but \emph{passive}---they only noticed
what a worker happened to read. So the new version's key addition is the
\textbf{active probe}. Each check can carry a small, deterministic re-read of
the world (a HEAD request, a schema fingerprint, a status re-check) that runs on
its own, with no LLM involved. Three things, each forced by what we measured
(\S6), shape how probes run:

\begin{itemize}
  \item They run on a \emph{separate channel} that does not disturb the world
  they measure---replay found three concrete ways a naive probe perturbs its
  own measurement.
  \item They fire \emph{on events and risk}, not on a fixed clock---a naive
  fixed schedule once cost 2,322 probe calls on a 107-call run.
  \item A non-obvious alarm interrupts \emph{only with a confirming probe}. The
  first version's ``wait for a second signal'' rule failed because correlated
  noise confirms itself---6 of 18 false alarms passed it. (A clear status-code
  error still interrupts on its own; no false alarm ever carried a status
  $\ge$ 400.)
\end{itemize}

Probes that compare a value need something clean to compare against. So at the
very start of a run, before any worker acts, the system takes one clean reading
of every surface it will watch, and later compares against that (deviation D30).

\textbf{We removed the judge.} The first version put an LLM ``judge'' between an
alarm and an interrupt. Its record: it approved all 34 alarms it ever saw,
ruled the only two genuine signals it down-ruled as noise, and once approved a
hallucinated misreading at 0.98 confidence. We committed---in writing,
\emph{before} the diagnosis results were readable---to making the no-judge
version the primary design (deviation D33 keeps a rebuilt judge only as an
exploratory side arm). There is no after-the-fact selection here.

\textbf{The check-writer is never told the fault list.} The first version
walked a list of eight named fault categories. The new version's check-writer
gets \emph{no list}. It reasons from two general ideas instead: what does this
step trust about the world, and which of six general \emph{shapes of change}
might break it (a field vanished, a status moved, a structure changed, a value
moved, an order scrambled, a relationship broke). Its only worked examples come
from the five fault types we had already seen, chosen by a rule fixed before any
prompt tuning (deviation D27). We attach category labels \emph{after} detection.
This is what makes the held-out test honest: the writer has to \emph{transfer}
general reasoning to a sealed fault type, not look it up (\S7). The pilot's
sharpest lesson is why this matters. Fault types differ enormously in how
visible they are. The one type the first version caught well was the loud,
repeating, status-coded one; the quiet, seen-once types were nearly invisible.
The probe exists to make every type as visible as the loud one.

\section{TripwireBench}
To test a monitor, you need a world that breaks on cue.

TripwireBench is that world: a fully mocked, deterministic set of services
(auth, inventory, pricing, shipping, a document store, a code repo), with
seed-derived data and no clock values anywhere. Faults are exact state changes
that fire at exact tool-call counts, so any run can be replayed byte-for-byte
from its setup. Because a fault fires at the same \emph{point in the work}
across systems and seeds, batch and treatment runs are directly comparable.
Four task types (API aggregation, repo migration, research synthesis, document
pipeline) run on a real agent CLI; one piece of middleware counts calls, fires
faults, checks tripwires, and writes a full trace.

\textbf{Every fault has to actually hurt.} No fault enters the experiment until
it provably breaks the plain batch baseline: injected runs must fail while clean
runs pass, under a rule fixed in advance. The pilot qualified all 9 fault--task
pairs. Two more fault types were authored \emph{after} the pilot and held out: a
metered quota that only complains after the budget is already spent
(RESOURCE\_BUDGET), and a silent version bump that quietly truncates a default
page of results (DEPENDENCY\_VERSION). Their exact parameters were sealed with
someone outside the project, so part of the scoring could not be over-fit. Both
cleared the same bar---the quota fault on the first try, the version fault only
on its third revision, with the two failed designs reported as findings rather
than hidden. The qualification record is itself instructive: a competent batch
planner kept \emph{defeating} faults---by paginating defensively before it saw
anything, or by reading collections once at the start so a later change landed
after the only look---and \emph{strengthening} the baseline un-qualified two
faults that had qualified earlier, so wounds must be re-checked whenever the
baseline improves (deviation D23).

\textbf{The systems we compare.} S1, plain batch (dispatch, wait, aggregate,
one redo). S2, naive interrupt (workers raise any anomaly straight to the
orchestrator---no check-writer, no judge). S3, a cost-matched periodic
re-check, tuned so its overhead equals our monitor's---the strongest fair
version of the obvious alternative. S4, our monitor without the judge. S5, our
full first version. The confirmatory study runs four arms---S1, S2, S3, and the
redesigned monitor---setting the unbuilt rebuilt-judge arm aside (deviation
D33).

\textbf{What we measure.} Wasted work (tokens and dollars spent after a fault
that don't reach the final answer); detection speed (time to detection, TTD:
calls from fault to first pause); detection rate; false-alarm rate (false-interrupt
rate, FIR: with no fault and no alarm scored as zero); extra cost on clean runs;
and plain pass/fail task success---no LLM grading anywhere. Pause attribution follows a written rule, with a hand-audited
sample (13 of 13 agreed).

\textbf{Integrity machinery.} The benchmark's credibility tooling is part of
the artifact: a frozen pre-registration with cryptographic fingerprints, a
numbered log of every departure, fault triggers that only the test harness can
reach, a log of every run touching any treatment component, a data embargo
before the bars were frozen, and sealed held-out parameters verifiable against a
public hash.

\section{The Pre-Registered Pilot and Its Verdict}
\textbf{Setup.} 13 task variants (9 qualified injected pairs + 4 clean) $\times$
5 systems $\times$ 3 seeds = 195 runs, all banked, with no malformed traces.
Tooling was pinned, model assignments fixed per role, sessions isolated. Faults
fire at the halfway point of each task's typical length, identical across
systems and seeds. Before the main experiment, two independent outside models
scored the written checks against a four-part rubric (95.1\% agreement) and they
passed at 89\% coverage against an 80\% bar. The four pass/fail bars---the kill gates, KG1--KG4---were frozen
before the first run, each with its consequence decided in advance. The decision
window after the verdict was itself pre-set: 48 hours, ``do not negotiate with
the data''.

\textbf{The verdict (computed once, shown verbatim):}

\begin{lstlisting}
KG1 recall: 35% (>=60%) | categories >=50%: 1/5 (>=4)            -> FAIL
KG2 FIR: S5=1.0 S2=0.0 S4=1.0 (S5<=0.5*S2 and S5<=0.7*S4)        -> FAIL
KG3 cost: S5 med $1.178952 vs S1 med $0.340831; success S5=4%
     S1=22%; clean overhead OVER (<=12%)                          -> FAIL
KG4 vs heartbeat: wasted med S5=15014.0 S3=7950 (>=20% better)
     | TTD med S5=3 S3=9 (>=2x)                                   -> PASS (TTD arm)
\end{lstlisting}

Plainly: the bars we froze returned \emph{fail} on our own flagship, and the
pre-decided consequences ran as written. Detection at 35\% fell in the kill
band (below 40\%), so that claim died. The false-alarm gate failed, so we
dropped the judge. The cost gate failed, which triggered a break-even fit that
later killed the cost claim outright (\S8). Only one thing survived: speed. Our
checks caught the fault in 3 calls to the periodic re-check's 9. But the
wasted-work side of that same gate failed (15,014 vs.\ 7,950 tokens):
\textbf{speed did not turn into savings under the first version's machinery}.

\textbf{Reading it carefully.} ``Detection'' here means a confirmed,
trace-attributable pause, defined before the bars computed. A more generous
reading (counting alarms that didn't pause) recovers only two runs---so the low
number is real silence, not a strict definition. By category: the one passing
type (PERMISSION\_AUTH, 5/6) is exactly the loud, repeating, status-coded one;
the quiet types scored 2/6, 1/6, 1/6, and 0/3. The predicted near-zeros showed
up.

\textbf{The naive baseline, stated plainly and prominently.} S2---workers raise
any anomaly straight to the orchestrator, no check-writer, no judge---caught
\textbf{15 of 27 (56\%) at a false-alarm rate of zero}, and beat or tied our
full first version in \emph{every} category. Workers given no filter escalate
rarely, and mostly correctly, at this scale. The folk worry that naive
escalation means alarm spam---which our false-alarm gate had baked in---was
simply false here. Sophistication actively hurt. S2 is now a required
head-to-head arm, under a rule we committed to in advance: if the simple thing
wins, we say so in the results, not a footnote.

\textbf{It hurt healthy runs, too.} On clean runs with nothing wrong, the full
first version succeeded 3 of 12 times against batch's 9 of 12. Seven of the nine
clean failures were self-inflicted: five runs ground to death adjudicating
noise, and two replanned so early that there was nothing to carry forward. All
34 interrupts had been judge-approved; 18 matched no fault at all, and 3 fired
on entirely clean runs. An uncalibrated monitor did not just fail to help. It
did harm.

\section{Why It Failed}
\textbf{Where the misses happened.} We classified all 27 injected treatment
runs along the chain \emph{write check $\rightarrow$ fault fires $\rightarrow$
alarm $\rightarrow$ judge $\rightarrow$ pause $\rightarrow$ attribute}. Of 18
misses: 1 had no check written; \textbf{12 had a check that simply never
fired---the dominant case}; 2 fired but never raised an alarm; 2 were ruled
noise; 1 fault never fired. So the problem was not blindness at write
time---coverage was near-complete once a measurement artifact was fixed. It was
\emph{silence at look time}.

\textbf{We did not grade our own homework.} We sent the diagnosis and the
redesign to two outside red teams for blind review, then settled every testable
claim with a deterministic replay over the recorded runs---with our commitments
fixed before any replay result was readable, and at zero LLM cost. All 27
injected worlds replay identically, which is what makes this possible. The
reviewers split roughly evenly on the eighteen claims they raised, which is
exactly why we \emph{adjudicated} their review with evidence instead of
deferring to it. The decisive checks: the story that evidence reached the system
and was lost is refuted (9 of 12 starved runs had \emph{no} post-fault reading
at all---nothing to lose); the story that false alarms replanned the surface
away is partial (that exact signature appears 0 of 12 times, though 7 of 12 died
of noise); and the story that runs simply died before reaching the fault is
refuted (0 of 17).

\textbf{The load-bearing result.} Our first diagnosis blamed the world: some
surfaces are only ever read once, so nobody could re-look. The same-seed
analysis corrected that, and it is the strongest evidence for the ``not about
intelligence'' half of the finding.

On identical seeds, the ordinary baseline systems re-looked at every starved
surface---12 of 12---and our ``smarter'' monitor did not.

The dumb baselines saw what the smart monitor missed. What starved it was its
own noise: every false alarm spends turns and budget that detection itself
needed. One run logged 172 alarms on healthy surfaces before grinding to death.
One mechanism, four symptoms: noise starved detection, the over-trusting judge
made the noise, and the grinding wrecked clean runs and cost. The opening story
(\S1) is this in a single run.

\textbf{Errata, kept visible.} Two early claims were corrected by the replay and
stay in the record with their corrections: the ``read-once surface'' story was
wrong for its own example (the surface was read \emph{after} the fault, so the
original telling was backwards), and ``one matcher defect'' became a \emph{class}
of dead checks---84 checks that could never match any real path under the live
dialect, fixed with a compile-time liveness test (deviation D24).

\textbf{We tested the fix before building it.} We mechanically built probes for
all nine pilot faults and replayed them against the broken worlds. Probe-first
confirmation blocked \textbf{all 18 false alarms}, and projected detection
\emph{rose} to 21 of 27, because probes recover what alarms lose. We are careful
here: \textbf{every number in this paragraph is a ceiling under stated
assumptions}---hand-built probes the real check-writer had not yet produced,
fixed run paths, the same nine known faults. It is a feasibility test of the
mechanism, and we quote it nowhere as a result.

\section{The Redesign, Measured}
The confirmatory study is the rerun path we wrote down \emph{before} the pilot
verdict existed: report the first version as it stands, diagnose from traces,
change the design with cited evidence, and re-measure against fresh bars.

The integrity rules are in ink. Every shared threshold is inherited from the
\emph{pre-verdict} values---60\% detection, the same category and floor rules,
the 2$\times$ speed bar---so no bar was set by anything the diagnosis told us.
New rules add a hard ``probes must be valid'' gate, a recovery-quality clause
(detect-only, detect-and-recover, and detect-and-justified-stop are counted
separately, so stopping cannot be dressed up as success), and \emph{absolute}
noise caps in place of the broken ratio gate. Fault settings were redrawn from
fresh seeds the designers never saw. The two sealed fault types sit inside the
scoring. A data embargo and a full run log stood throughout. And because the
check-writer never sees the fault list, the sealed types are unseen even in the
vocabulary it reasons with.

\textbf{Result.} The experiment ran to completion---\textbf{172 of 172 runs
banked}, across four systems (S1, S2, S3, and the redesign)---and the bars
computed once.

\begin{lstlisting}
1bKG1 detection: recall 10/15 = 66.7% (>=60%, floor clear);
     5/5 seen categories >=50% (API 6/6, SCHEMA_DRIFT 3/3, AUTH 6/6,
     TOOL_CONTRACT 3/3, RETRIEVAL_INTEGRITY 3/3); probe-validity 10/10;
     recovery-quality 17/24 = 70.8%             -> detection sub-terms PASS
1bKG2 noise/self-harm: clean median FIR 0, P95 0, max 0; 0 grinds;
     pre-detection median FIR 0; clean success 8/12 = 66.7%        -> PASS
1bKG3 cost: V2 clean median $0.3642 vs batch $0.2342;
     clean overhead 55.5% (cap <=12%)                             -> FAIL
1bKG4 vs heartbeat: V2 wasted 7008 vs S3 6404 = 1.09x;
     waste-parity not met (S3 detects 0/31)                       -> FAIL
OVERALL                                                           -> FAIL
\end{lstlisting}

The verdict is \textbf{fail, driven by cost}---strong detection that misses the
cost bar.

One reporting note we keep visible. An admissibility check (a byte-identity
replay) came back 123 of 124 identical; the single difference is a harmless
ordering quirk on the non-detecting periodic-re-check arm, characterized and
logged as deviation D34. The gate \emph{code} folds that check into the
detection gate, so the combined gate prints \texttt{FAIL} even though all four
detection sub-terms pass. We report it split out---detection passes; the
admissibility check is fine bar one benign diff---and note the overall fail is
unaffected, because it rests on cost.

\textbf{The 12\% cost cap was a judgment call, and it is not load-bearing.} We fixed
the clean-overhead cap at 12\% in the pre-registration, before any data existed. It was
a pre-committed judgment, not a quantity derived from a model.
The frozen pre-registration records no rationale for the number---it appears only
as the KG3 clean-overhead threshold. It was a pre-committed authors' judgment,
disclosed as such. Nothing
downstream turns on the exact number. The verdict is invariant to the cap's value: the
redesign fails any cap below 55.5\%, and the second-family floor of $+17\%$ (\S9) fails
any cap below 17\%---every bar a deployer could plausibly have set returns the same
verdict. Post-hoc, the cap is if anything \emph{generous} to the monitor: located on the
fitted break-even curve (Appendix~B), a 12\% clean overhead pays for itself only near a
41\% fault rate even under an accounting that ignores replan cost, and at plausible fault
rates (5--25\%) the overhead that would break even is 1.5--7.3\% (all post-hoc).

\textbf{Why the redesign wastes more than the do-nothing heartbeat (a post-hoc autopsy).}
The redesign's median wasted work (7,008 tokens) exceeds the cost-matched heartbeat's
(6,404) even though the heartbeat detects nothing. An additive decomposition of the mean
gap (807 tokens) puts the entire excess in one bucket: re-dispatch rework. When the
redesign detects a fault it replans and re-dispatches workers, discarding the partial
output of the workers left on the doomed plan; that discarded rework averages about 1,055
tokens per non-clean run, appears in 8 of 31 cells, and the heartbeat---detecting nothing,
replanning nothing---discards none of it. It is not sunk pre-fault cost: the redesign
actually books \emph{less} of that than the heartbeat (mean 2,278 vs.\ 2,806 tokens),
because it stops early and truncates the exposure window the heartbeat runs to completion.
So the 1.09$\times$ parity is the price of corrective action: the redesign pays to detect,
stop, and re-plan around the fault, and the metric counts that recovery machinery as waste,
while the heartbeat's cheaper run simply finishes on the broken plan and delivers the wrong
answer. Excluding the re-dispatch rework, the redesign's median waste falls to 6,117 tokens,
below the heartbeat's 6,404. The frozen 1.09$\times$ verdict line stands; the autopsy shows
the gap is recovery cost, not the monitor wasting more.

\textbf{Detection, in plain numbers.} The redesign caught \textbf{24 of 31}
injected runs (Wilson 95\% lower bound 60.2\%). The passive baselines caught
\textbf{0 of 31} (both batch and the periodic re-check), and the naive baseline
caught \textbf{12 of 31}; the redesign's edge over that naive baseline is
significant (one-sided Fisher exact $p=0.002$). On the five seen fault types it
was perfect, and it finally caught the quiet, seen-once types the first version
could not (\S5). Two numbers describe the same runs from different angles: 24 of
31 \emph{detected} (77\%), and 10 of 15 (66.7\%) on the gated subset---the subset
restricted to cells that admit a recovery path, a rule frozen in the
Phase-1b pre-registration (prereg\_1b.md, commit \texttt{6c8cc47}, 2026-06-12), which
labels each injected cell recoverable iff it admits a recovery path. They differ by \emph{definition}, not contradiction; and the
choice of subset is not outcome-determining, because 24 of 31 (77\%) also clears
the 60\% detection bar (24 of 31 $=$ 77.4\%). The 24 detections split into 3 detect-and-recover, 14
detect-and-justified-stop, and 7 detect-only, with 17 of 24 (70.8\%) in the two
passing buckets. The five seen categories are each perfect, but three sit at
n=3, where the Wilson 95\% lower bound is only about 44\% and one run swings the
rate 33 points---always reported with the bound.

\textbf{It transferred, directionally, to a fault type it had never seen.} The
redesign caught the sealed RESOURCE\_BUDGET type \textbf{3 of 5} times (Wilson
95\% lower bound 23\%), where every passive baseline caught it \textbf{0 of 5}.
With $n=5$ the one-sided Fisher exact test does not reach significance
($p=0.08$), so we read this as a \emph{directional} transfer of general reasoning
to a fault type held out of its training entirely---not an established effect.

The second sealed type, DEPENDENCY\_VERSION (a silent version bump that
truncates a default page), was missed by \emph{every} system, ours included
(\textbf{0 of 5}). Read-only trace checking over all 20 of those runs explains
why. The fault fires on the very first tool call---\emph{before} the redesign
takes its clean opening reading---so even our reference reading is already
truncated, and the later re-reads match it. And no surface anywhere carries a
version field, so the passive and naive baselines have no error to flag either.
Because \emph{every} system missed it---including the anomaly-raising naive
one---this is the observation bound of \S6 meeting the design on the read side,
not a wart specific to our monitor. Stated as a scope condition over \emph{all}
the detections: a reference-comparison monitor catches a fault only if it lands
\emph{after} the monitor's opening reading and \emph{before} its last look. The
benchmark's main faults fire at the halfway point---inside that window by
construction---which is why they are catchable at all; DEPENDENCY\_VERSION fires
on the first tool call, before the window opens, which is why it walls every
comparison-based system, ours included.

\textbf{The first version's pathology, reversed.} The first version's
signature---a rubber-stamped flood of false alarms that starved detection and
hurt clean runs---is gone. Clean-run false-alarm rate is 0 (median, 95th
percentile, and max) in the noiseless world (no benignly volatile fields; \S10),
with zero grind-deaths. Its clean-run success (8 of 12)
trails batch's (9 of 12) by a single run, and all four of its clean failures
occurred at \emph{zero} interrupts---worker-deliverable shortfalls (missing
citations, unset fields, validation not run) that the no-monitor baseline shows
too, so the monitor neither caused nor could have prevented them. And the win
holds under real
pressure: the clean set includes the legitimate config-migration write that used
to trip a false alarm, and the redesign ran it at zero interrupts---while the
naive baseline raised false alarms on four clean runs. The redesign bought back
detection \emph{and} dropped noise to the naive baseline's zero floor. Against
the naive baseline's 12 of 31, sophistication helped this time. But the help did
not survive the cost gate---which is the subject of the next section.

\section{The Coupling: One Result, Not Three Failures}
\textbf{Start with the prediction economics makes.} A cost paid \emph{once per
run}---read the plan, write the checks---should amortize across a large fan-out
and all but vanish per worker: one fixed charge spread over many parallel
workers. That is the natural reason to expect plan-compiled monitoring to pay for
itself at scale. The measurement returned the opposite direction, and this
section is why.

The cost failures of the first and second versions are not two missed targets.
They are one fact, seen from a few sides: \textbf{detection and the cost of
restoring observation are tied together}, because the one thing that restores
observation is also the main cost.

\textbf{The cost is the capability.} We broke the 55.5\% clean overhead down
call by call. Almost all of it is one new component: the once-per-run LLM step
that reads the plan and writes the checks (median \$0.1376 on clean runs,
against a \$0.130 gap over batch), present in all 12 clean runs. The whole
detection machinery---the probes, the re-readings, the confirmation
logic---costs essentially nothing measured. So the overhead lives in the one
step we cannot remove without removing the capability: that step \emph{is} the
reasoning that produces both the detection and the transfer to unseen fault
types (\S7). (\S9 sharpens this: on a second model family the load-bearing content
of that step turns out to be its \emph{coverage} of the plan's assumptions, not
its reasoning effort.) Making it cheaper---caching it, shrinking it, skipping
it---is not a free optimization. It risks the detection it pays for. That mutual dependence is
the coupling. A no-check-writer variant makes the point from the other side:
drop that step and the system is nearly cost-neutral with batch on clean
runs---so the step really is the cost center, though that variant's detection is
only established on its narrow home case, so we report it as cost isolation, not
a free lunch.

\textbf{The obvious fix doesn't open.} A fixed per-run cost should shrink, per
worker, as you add more workers. So the natural escape is to spread that
one-time check-writing cost across a bigger fan-out. The break-even model even
predicts a crossover for the redesign---but that is a model prediction, not a
measurement (Appendix B).

We built a width-scaling benchmark and ran a sizing pilot to test it. The
pilot's point estimate---kept out of the formal sizing by design---is that the
redesign is \textbf{costlier than batch at every width and fault rate we tried,
and the gap \emph{grows} as we add workers}, because the check-writing scales
with the planning the system does, not with the parallel work it could dilute
against.

We treat this as a \emph{closing negative}, not a confirmation: the escape shows
no sign of opening, and nailing it down at the pre-registered precision would
cost about ten times the study budget ($\approx$ \$4,430 against a \$450 budget)
for a number already pointing the wrong way---so a full fan-out study is not
warranted. The price tag is the side note. The \emph{direction} is the finding.

\textbf{That closes the width escape; a depth escape we do not close.} The sizing
pilot varies \emph{width}, the number of parallel workers. A second escape is
\emph{depth}: on longer tasks, execution tokens might grow faster than the plan,
so the fixed check-writing cost could amortize against a longer run rather than a
wider one. Our benchmark does not vary task depth enough to settle this---across
the four task types execution size is nearly constant while plan size varies
threefold and largely independently of it (Pearson 0.30, post-hoc)---so we do not
claim the depth escape is closed. What the data does show is the scope condition:
the plan is a non-trivial fraction of execution in every task type (0.32 to 0.86
of worker tokens). We therefore scope the coupling to runs whose plan is a
non-trivial fraction of execution---which is precisely the multi-agent fan-out
regime this paper targets.

\textbf{And it fails even when looking is free.} Our benchmark is a
deterministic mock, where re-reading the world costs nothing. That is the
\emph{most generous} setting the system could ask for---it charges nothing for
the very looking the system exists to do. And it \emph{still} misses the budget,
on the one-time check-writing step alone. So the measured 55.5\% is a
\emph{floor}: in a real deployment, probes hit real endpoints at real cost, and
the volume that reads as free here comes back as real overhead. A system that
cannot pay for itself where looking is free will not pay for itself where
looking costs money.

\textbf{Putting it together.} Restoring observation fixes detection (\S7) and
ends the false-alarm self-harm (\S7). But the only place to cut cost is the
capability-critical check-writing step, the obvious amortization does not open,
and the failure holds even with looking priced at zero. So in budget-bounded
multi-agent execution, \textbf{detection and cost are coupled, and this kind of
monitoring does not yet pay for itself.}

\textbf{The coupling has one root, on both model families.} You cannot monitor a
plan without reading the plan, and reading the plan is the irreducible cost. On
the incumbent writer that cost shows up as a reasoning step; on GPT-5.5 (\S9)
reasoning is cheap and detection-free, and the floor is instead the writer's
input tokens---the plan itself---spent through checklist coverage. The claim is
not that thinking is expensive. It is that \emph{coverage of the plan's
assumptions has a floor cost proportional to the plan}, and coverage is what does
the catching.

\textbf{The verdict is conditional on how value is priced.} The coupling is a
claim about the \emph{cost} side and holds under any value model. ``Does not yet
pay for itself'' is a claim relative to a \emph{token-denominated} value
model---the one every bar in this paper uses. But the 14 detect-and-justified-stop
runs each prevented a confidently-wrong deliverable---the exact outcome of this
paper's opening run---and every bar here prices that prevented error at zero.
Under a correctness-denominated value model the economics are open and
unquantified here; pricing that term is the next experiment. As a post-hoc frontier
(not a deployment claim), the arithmetic can at least be stated: at clean overhead
\$0.13 per run and justified stops on 14 of 31 injected runs, one prevented wrong
deliverable breaks even near $\$0.13 \times 31/(14p)$ in the same token units---about
\$2.90 at a 10\% fault rate and \$1.15 at 25\%---complementing Appendix~B's Table~3,
which prices the same question from the overhead side. The bars measure the
token refund, and no one buys insurance for the refund. That next problem---and it
is a fresh experiment, not a patch---is what the coupling leaves open.

\section{The Coupling Holds on a Second Model Family (Live)}
Everything above is one model family: a Claude check-writer, Claude and Haiku
workers. A fair objection is that the coupling is that family's artifact. So we
tested it on a second one.

\textbf{This study was pre-registered separately, and later.} After the original
pre-registration and both its verdicts, on 2026-07-01, we froze a new, standalone
pre-registration for this study alone---its own 12\% bar, its own arms, its own
pre-decided consequences, its own hard spend cap, all committed before any of its
data existed. It is not part of the original plan and we do not present it as one:
it is an additional, later-dated extension that reruns the paper's method on a
different vendor's model. Folding it into the original frozen pre-registration
would break the exact discipline this paper is about, so we keep the two apart
and dated.

We swap GPT-5.5 (pinned \texttt{gpt-5.5-2026-04-23}; \$5.00 / \$0.50 / \$30.00 per
million input / cached-input / output tokens, verified 2026-07-01 \cite{ref38})
into the check-writing step through the same frozen seam the redesign already
exposes. Everything downstream---grounding, probes, confirmation, scoring---is
byte-identical; only the writer changes.

\textbf{Offline, the escape looked open.} GPT-5.5 is a cheaper and cleaner writer
than the incumbent. It writes valid checklists at \$0.115 per call against the
Claude writer's \$0.1376 (\S8), and it emits the value-comparing check that the
cheap alternative never does (median 8 such checks per checklist, against 0 for a
budget model). We first scored its checklists in the zero-LLM deterministic
replay of \S6. There, detection is \emph{flat} as we turn its reasoning effort
down---\texttt{none}, \texttt{low}, and \texttt{medium} all catch the same 16 of
18 seen faults---because the reasoning tokens are not what does the detecting. At
\texttt{effort=none} it costs \$0.064 a call, well under the incumbent, detection
intact. On the reproducible replay, the coupling looked beaten.

\textbf{Live, it was not.} The replay prices looking at zero and holds the world
byte-identical; it cannot see what a shorter checklist costs in a real run. So we
ran the priced matrix live---real workers, real timing---and measured the
overhead the bar is written on. The cost lever that could actually clear 12\% is
not reasoning effort (which is detection-free) but \emph{checklist coverage}: a
shorter checklist is a cheaper writer call. We capped it with a single general
instruction, committed in the pre-registration and blind to the fault list---keep
the $N$ most load-bearing assumptions, judged only by how badly the plan breaks
if each is false, never by what kind of fault it represents---at three frozen
settings (Table~\ref{tab:v3-overhead}).

\begin{table}[t]
\caption{Live clean-run overhead by checklist-coverage cap (GPT-5.5 writer,
reasoning fixed at \texttt{low}). Batch (S1) clean median \$0.2219 is the
denominator, measured live. Single seed $\times$ 4 tasks (departure D-V3-1); the
writer/batch-worker column is the robust reading (see text).}
\label{tab:v3-overhead}
\begin{tabular}{lrrrr}
\toprule
coverage & assumptions & writer \$ & overhead & overhead \\
 & (median) & (median) & (frozen) & (writer/batch) \\
\midrule
full   & 40 & 0.093 & $+33.9\%$ & $+42.0\%$ \\
cap-12 & 12 & 0.049 & $+20.6\%$ & $+21.9\%$ \\
cap-6  &  6 & 0.038 & $+22.4\%$ & $+17.1\%$ \\
\bottomrule
\end{tabular}
\end{table}

\textbf{No setting clears 12\%.} The raw
$(\mathrm{monitored}-\mathrm{batch})/\mathrm{batch}$ median is noisy at this
study's single seed---worker cost varies run to run and swamps the writer
saving, so cap-6 even reads worse than cap-12 on that column---so we read the
overhead the robust way, as the writer call against the batch worker cost it adds
to: $+42\%$, then $+22\%$, then $+17\%$, every setting above the bar. The floor is
the writer's own input tokens; even a six-assumption checklist costs \$0.038.

And the cap that lowers cost costs detection---and, extended to three seeds on
the injected cells, \emph{destabilizes} it (Table~\ref{tab:v3-detection}).

\begin{table}[t]
\caption{Live detection across three seeds (injected cells): the number of seeds,
out of three, in which the fault's interrupt fired during full worker execution,
by checklist-coverage cap. Five catchable faults plus the DEPENDENCY\_VERSION
wall. The bottom rows give catchable-fired-of-five per seed, then the median and
range.}
\label{tab:v3-detection}
\begin{tabular}{lccc}
\toprule
injected fault & full & cap-12 & cap-6 \\
\midrule
removed endpoint (API\_SURFACE)        & 3/3 & 3/3 & \textbf{0/3} \\
expired token (PERMISSION\_AUTH)       & 3/3 & 3/3 & 3/3 \\
skipped gate (TOOL\_CONTRACT)          & 3/3 & 3/3 & 2/3 \\
swapped passage (RETRIEVAL\_INTEGRITY) & \textbf{2/3} & \textbf{1/3} & \textbf{1/3} \\
RESOURCE\_BUDGET (quota cliff)         & 3/3 & 3/3 & 2/3 \\
DEPENDENCY\_VERSION (silent bump)      & 0/3 & 0/3 & 0/3 \\
\midrule
catchable-of-5, per seed & 5, 4, 5 & 4, 5, 4 & 3, 1, 4 \\
\quad median (range)     & 5 (4--5) & 4 (4--5) & 3 (1--4) \\
\bottomrule
\end{tabular}
\end{table}

\textbf{The direction is robust; the staircase is not.} The load-bearing claim
holds on every seed: the short checklist catches less than the full one---5, 4,
and 5 of the five catchable faults at full, against 3, 1, and 4 at cap-6, so full
beats cap-6 on all three draws. At the median the slope is a clean 5, 4, 3. But
the per-rung counts are seed-variable, and prominently so: cap-6 ranges from 1 to
4 of 5, and one seed is non-monotone (4 at full, 5 at cap-12, 1 at cap-6). The
tidy staircase is an artifact of a single draw; the \emph{descent} is not.

\textbf{Cheapening the monitor does not merely lower detection---it destabilizes
it.} At full coverage detection is fairly reliable (4--5 of 5); by cap-6 it is
timing-dependent (1--4 of 5). The cap does not walk a steady dial down; it turns a
dependable monitor into a gamble on whether six probes happen to land where and
when the fault fires. The instability is itself a cost of cheapening, and it is a
finding, not noise.

\textbf{And the fragility sits where the theory puts it.} The value-shaped fault
(a swapped passage) is the one needing a probe on a \emph{specific} surface at the
\emph{right} time---the closest fault to the observation bound of \S6. It is
marginal even at full coverage (2 of 3 seeds), and 1 of 3 at both caps: the
hardest-to-see fault is the first to wobble and the first to fall. That is the
coupling's own logic surfacing as variance, not a wart in it. By contrast the
removed endpoint is robust until it is not (3 of 3 at full and cap-12, then 0 of 3
at cap-6, where six assumptions no longer reach the pricing surface at all); the
auth fault, which 401s every surface at once, fires 3 of 3 everywhere; and the
version bump stays the \S7 wall (0 of 3 at every setting, excluded from the
slope). The one setting that catches all five on a good draw is the full
checklist---which has the \emph{worst} overhead. Cheapening the writer loses
faults, makes the rest unreliable, and \emph{still} misses the bar.

\textbf{These three draws are not a controlled seed replication, and we do not
claim one.} The dominant source of the run-to-run variation is the \emph{live}
non-determinism of the unseeded orchestrator, workers, and writer; the world seed
adds fixture and token-stream variation on top. That is the right variation for
this question---it is exactly what made the live numbers timing-variant---so the
draws test robustness to deployment reality rather than to a re-seeded fixture. It
is not a claim of three independent controlled trials.

\textbf{Verdict, against the frozen V3 bar: the coupling holds (live).} Both
committed failure conditions are met---detection erodes as cost drops, and no
setting clears 12\%. This is the paper's thesis on a second model family, and it
arrives through a \emph{new mechanism}. On our own system the cost center was the
one reasoning step (\S8); on GPT-5.5 the reasoning is cheap and detection-free,
and the cost lever is instead the checklist's \emph{length}---which is its
coverage. Different mechanism, same coupling: the thing you would shorten to save
money is the thing that does the catching. The priced matrix and the multi-seed
extension together cost \$6.54 in metered spend, under the pre-registered cap.

\textbf{What this is and is not.} It is a fail on the bar, reported as one, not an
escape. GPT-5.5 is genuinely the better writer; it does not clear the cost. Four
cautions survive verbatim. (1)~The clean cost cells are a single seed (departure
D-V3-1); the injected detection cells were extended to three seeds. The robust
signals are the writer-cost overhead (clean) and the detection \emph{direction}
(full beats cap-6 on all three seeds); the per-rung counts are seed-variable, and
we report that instability above as a finding rather than bury it as noise.
(2)~Prompt-caching did not engage this run; a warm cache could push the
six-assumption cost floor toward 6\%, but that was not achieved here and is future
work, not a result---and even a cost pass would not matter, because at cap-6
detection already loses faults and is unstable (1--4 of 5), so the detection-axis
coupling stands regardless. (3)~The offline replay and the live run are different
measurements and we do not blend them; the offline optimism was contradicted live
at the load-bearing point, and that contradiction is part of the finding, not an
embarrassment. (4)~The DEPENDENCY\_VERSION wall of \S7 stands for GPT-5.5 too---
missed at every setting, every seed, constant across arms, and so excluded from
the coverage slope; and RESOURCE\_BUDGET, which reads 8 of 8 in the coverage-bound
offline replay, fires live at full and cap-12 but only 2 of 3 seeds at cap-6---
seed-fragile once the checklist is shortest. Each is labeled as what it is.

\section{Threats to Validity}
\textbf{These are findings, not laws.} One harness, one mocked world, three
seeds, bars written on medians. The workers are one model family throughout; the
\emph{check-writer}---the coupling's cost center---we now test on two (\S9), where
the coupling reproduces through a different mechanism. Plan shape varies run to
run and we do not control it. Our injected faults approximate, but do not sample, the real
distribution; qualification guarantees each one breaks the baseline, not that it
is representative. We state the observation and cost results as regularities of
this setting, not proofs.

\textbf{The free-looking mock cuts both ways.} It makes the cost verdict
\emph{conservative}---the system fails the budget where looking is free
(\S8)---but it cannot settle real-deployment cost, where probe volume and
latency return. The 55.5\% is a floor, and the one cost lever is also the
capability, so the cost fix carries detection risk. Both are named; neither is
solved. The 12\% bar itself was a disclosed pre-registered judgment, not a
derived quantity, and the cost verdict does not turn on it: every cap a deployer
could plausibly set returns the same fail (\S7).

\textbf{The mock generates no benign noise, and that bounds the other verdict.}
The same free-looking mock that makes the cost verdict conservative also emits
\emph{no} benign anomalies---no transient 500s, no rate-limit jitter, no flaky
reads, no latency spikes. So the redesign's zero-false-alarm result---its clean
false-alarm rate of 0---is measured in the most benign environment possible and is
an \emph{upper bound} on deployment behavior, not a deployment number. (The naive
baseline is not even zero here: it already false-fired on four clean cells in the
noiseless world, \S7.) The cascade is honest, and we measured it rather than leaving it hypothetical:
under injected benign noise the \S6 self-starvation mechanism did not manifest at
this scale (zero grind-deaths), but the content channel does draw false interrupts
while one-shot status transients stay isolated, and sustained noise is untested---so
the confirmatory detection numbers inherit that same, now-measured and scoped,
qualification. And the probes are not side-effect-free in the field:
here a probe costs nothing and perturbs nothing, but a probe against a
\emph{metered} surface spends the very quota it checks, so the RESOURCE\_BUDGET
probe is potentially self-defeating live. The mock is generous to the noise
result exactly as it is harsh to the cost result, and this paper leans on both
directions: we accept the conservative cost verdict, so we flag---and, in the probe below,
measure---the optimistic noise one.

A pre-registered benign-noise probe (A7, with a family-closing follow-up A7b, one
seed per cell) injected three recoverable non-fault anomalies---a transient
500-then-success, an elevated-latency value, and a backward-compatible extra
field---into the redesign and the naive baseline on clean tasks, with a noiseless
control isolating the noise from run-to-run variance on the same CLI. Each ran
under its own dated mini pre-registration with predictions frozen before the run;
all three frozen predictions were refuted and are retained in the record. Benign
noise produces false interrupts in the redesign, selectively by channel. Against
content noise---a backward-compatible extra field---the redesign false-fires: its
compiled schema probe is a full-field-set fingerprint, so an added field changes
the fingerprint and the monitor reports a schema-shape drift even though every
required field and value is unchanged (the noiseless control interrupted zero of
nine runs; the same seeds interrupted once the field was present). That
sensitivity is the capability, not a defect in itself: comparison against a
reference is what detects at all, and separating benign drift from a real break
would need the plan-level reasoning whose cost is the subject of \S8. Against a
one-shot status transient---a 500 that heals on retry---the redesign does not fire
on any surface, because its probes read a perturbation-isolated side channel and a
single healed transient on the worker path is essentially never co-observed in
this design (and never, by construction, under our injection); the naive baseline
does escalate on the transient, but the orchestrator dismisses it. No run suffered
a grind-death, so the \S6 self-starvation mechanism did not manifest here. The
confirmatory clean false-alarm rate of 0 therefore remains the noiseless-world
figure and an upper bound: it survives a one-shot benign transient but not a
benign content addition, and sustained status noise is untested---the bound does
not extend to it. Three limitations qualify the reading: one seed per cell with
nondeterministic agents (a matched-seed, not byte-identical, control); the added
field the monitor trips on is our own always-present envelope, a realistic but
synthetic anomaly; and the probe ran on a newer CLI than the frozen confirmatory
pin, which the noiseless control mitigates but does not eliminate.

\textbf{The redesign learned from the pilot's traces.} That is a real leak, and
we mitigate it rather than hide it: fresh fault settings the designers never
saw, sealed held-out types (drawn by an unseeded cryptographic process with a
committed hash) inside the scoring, and a check-writer whose blindness to the
fault list is verifiable by inspection. The held-out types are held out in their
settings, not their \emph{kind}---same hands, same ontology---and the baseline
defeated two of three designs before one qualified, which shows authoring
genuinely wounding faults is itself hard. A held-out-\emph{ontology} arm and a
natural-failure arm remain the main generalization tests for future work.

\textbf{We built the benchmark, the system, and the bars.} The counterweight is
the apparatus: frozen pre-registration, pre-decided consequences, cryptographic
fingerprints, binding outside review with trace-level rulings, a public hash,
and negatives reported at full prominence---including a naive baseline beating
our own design (\S5). The strongest evidence it binds is behavioral: frozen bars
returned a kill verdict three times, across two model families---twice on our own
system and, under a separately and later frozen pre-registration, a third time on
a different vendor's model (\S9)---and the pre-decided consequences ran as written
every time.

\section{Conclusion}
We pre-registered a system we believed in, and the bars we froze killed most of
it. We diagnosed why from fully recorded runs, rebuilt the detector around the
diagnosis---and the bars killed the rebuild too, on cost. Because each verdict
came with its mechanism instead of a shrug, they compose into one finding.

In budget-bounded multi-agent LLM execution, catching a broken plan is limited
by observation, not by the monitor's cleverness: the system that re-looks,
detects, and the plain baselines that kept re-looking saw what the clever one
starved itself out of seeing. And because the one thing that restores
observation---re-reading the plan to write fresh checks---is also the main cost,
detection and cost are coupled. Restoring observation fixes detection and ends
the self-harm. It does not yet pay for itself, and the obvious way to make it
pay does not open, even when looking is free.

And it is not one model's artifact. A separately and later pre-registered live
study put a different vendor's model in the check-writer's seat---a cheaper,
better writer---and the coupling held: the bar failed a third time, now through
the checklist's length rather than its reasoning (\S9).

That coupling is a claim about cost, and it holds however value is priced. What
it leaves open is the other side of the ledger: every bar here is
token-denominated, and each of the fourteen justified stops prevented a
confidently-wrong deliverable---the outcome of the opening run---priced at zero
throughout. Under a correctness-denominated value model the economics are open,
and pricing that term is the next experiment.

That coupling, not any single missed bar, is the result. It is what a
plan-driven monitor has to solve next.

\section*{Data Availability}
A replication package---the benchmark world and harness, the frozen
pre-registration with fingerprints, all 195 pilot traces and the deterministic
replay battery, the 172-run confirmatory experiment and its report, the cost
autopsy, the fan-out sizing pilot, the qualification and review records, the
departure log, and the run log---is at [REPLICATION-PACKAGE-URL]. It also includes
the second-family study of \S9: its own, separately and later frozen
pre-registration (dated 2026-07-01), the live priced matrix and per-run cost and
detection records, and the offline replay scores. It further includes the
benign-noise family (\S10): the A7 and A7b mini pre-registrations with their
ratification commits, the per-cell run ledgers, and the archived trace zips
(pilot, confirmatory, and benign-noise) verified by a committed SHA-256 manifest.
The held-out parameter values
stay sealed with someone outside the project; their fingerprint is committed
publicly so the file cannot be swapped.

\appendix

\section{Integrity and Departure Log}
Load-bearing departures, in brief; full records in the departure log (Data
Availability). \textbf{D23} removed two conditionally-qualified runs from the
experiment (215 $\rightarrow$ 172 with four arms) and named the \emph{write-side
single-look} residual---an agent's own write can overwrite a change it never
saw. \textbf{D24} fixed the dead-check class (84 checks, instrument-level, with a
regression test). \textbf{D25/D27} fix the held-out scoring outside the project
and freeze the check-writer's example-selection rule before any tuning.
\textbf{D28--D32} record the confirmation, scheduling, opening-reading,
write-handling, and grounding rulings of the rebuild, each committed before its
code and each preserving identical behavior when the flag is off. \textbf{D33}
set the unbuilt rebuilt-judge arm aside. \textbf{D34} characterizes the lone
replay difference (123/124, benign, on the non-detecting arm); no number or the
overall fail changes.

\textbf{The second-family study (\S9) is pre-registered on its own, and later.}
Its frozen pre-registration is dated 2026-07-01---after the original and both its
verdicts---with its own 12\% bar, arms, coverage settings, verdict definitions,
and a hard spend cap, all committed before its data existed and never folded into
the original pre-registration. \textbf{D-V3-1} logs its execution departures: the
clean cost cells ran at a single seed (four tasks) with a reduced injected set,
executed in resumable foreground chunks after the background runner proved
unreliable---execution-only changes that leave the frozen bar, lever, verdict
meanings, and spend cap untouched. A subsequent robustness extension re-ran the
\emph{injected detection} cells at two further seeds (three in total), reported in
\S9; the direction of the detection slope is robust across all three while the
per-rung counts are seed-variable. The study spent \$6.54 of metered budget across
matrix and extension, under its \$12 cap; no cap was hit.

\textbf{The benign-noise family (\S10) is pre-registered on its own, and later
still.} Two dated mini pre-registrations---the A7 addendum ratified 2026-07-02
(checkpoint \texttt{7a3807fe}) and A7b ratified 2026-07-03---froze arms, seeds,
spend caps (\$15 and \$8), and per-arm predictions before each run. \textbf{D35}
corrected the naive baseline's noiseless clean false-alarm characterization;
\textbf{D36} froze the latency envelope; \textbf{D37} logged a CLI-version drift
that \textbf{D38} then resolved, when the A7b noiseless control overturned A7's
first reading: the false interrupts were noise-caused---a closed-world schema
fingerprint firing on a benign added field---not task-intrinsic. Spend was \$8.75
and \$4.96, each under its cap. All three frozen predictions were refuted and are
retained in the record.

\section{Fan-out Break-Even Model (a prediction, not a result)}
Cost-positivity holds when
$C + J + p\cdot R < p\cdot(W_{\mathrm{batch}} - W_{\mathrm{sent}})$. Fitting the
model to the first version's runs ($C$~=~\$0.322 check-writing, $J$~=~\$0.048
judge, $R$~=~\$0.257 per replan) gives a \emph{negative} waste gap of
$-$\$0.072 and \textbf{no crossover at any fan-out, with probability 1.00 over
1,000 resamples}. Re-fitting to the redesign gives a positive gap and an
extrapolated crossover near 86, 40, and 25 workers at fault rates 0.10, 0.25,
and 0.50. \textbf{These are model extrapolations---predictions to test, never
measured results.} The sizing pilot (\S8) tested the direction and closed it
negatively.

\textbf{Locating the 12\% cap on this curve (post-hoc).} Inverting the same model
for the maximum clean overhead a deployer could tolerate---rather than for a
crossover fan-out---does not admit a clean answer at the \emph{measured} fan-out.
At $n_0=3$ the redesign's per-replan cost ($R$~=~\$0.221) already exceeds its
per-run waste gap ($W_{\mathrm{batch}}-W_{\mathrm{sent}}$~=~\$0.068), so
$C+p\cdot R<p\cdot\Delta W$ fails for every fault rate and every overhead
including zero: a \emph{free} monitor is cost-negative at $n_0=3$, turning
positive only past a fan-out of about ten workers. The nearest meaningful locator
drops the replan term---asking only whether the clean overhead is recovered by
the expected saved waste, $M<p\cdot\Delta W$, which is generous to the
monitor---and puts the 12\% cap at a fault rate of $p^\star\approx41\%$
(Table~\ref{tab:breakeven-locate}). Even under that generous reading 12\% was, if
anything, generous to the monitor, and the measured 55.5\% is far outside
break-even either way. All figures here are post-hoc.

\begin{table}[t]
\caption{Post-hoc: maximum clean overhead at which redesign monitoring is
cost-positive, by fault rate $p$, under the generous waste-recovery reading
($M<p\cdot\Delta W$; $\Delta W=\$0.068$, baseline \$0.234). The pre-committed
model including replan cost admits no positive break-even at the measured fan-out
($n_0=3$).}
\label{tab:breakeven-locate}
\begin{tabular}{lrrrr}
\toprule
fault rate $p$     & 0.05 & 0.10 & 0.25 & 0.50 \\
\midrule
max clean overhead & 1.5\% & 2.9\% & 7.3\% & 14.6\% \\
\bottomrule
\end{tabular}
\end{table}

\begin{thebibliography}{99}
\bibitem{ref1} Execution monitoring survey (Dunlap et al.).
\bibitem{ref2} Davis-Mendelow et al. Assumption-based planning. AAAI 2013.
\bibitem{ref3} Bercher et al. Plan repair. ICAPS 2014.
\bibitem{ref4} Bauer, Leucker, Schallhart. Runtime verification for LTL and TLTL. ACM TOSEM.
\bibitem{ref5} Havelund, Rosu. Synthesizing monitors for safety properties. TACAS 2002.
\bibitem{ref6} Ferrando, Cardoso. RVPLAN: Runtime Verification of Assumptions in Automated Planning. ICAART 2022, Vol. 2, pp. 67--77. DOI 10.5220/0010776500003116.
\bibitem{ref7} Parallelized planning-acting. arXiv:2503.03505.
\bibitem{ref8} Weak-to-strong monitoring. arXiv:2508.19461.
\bibitem{ref9} AgentSpec. ICSE 2026. arXiv:2503.18666.
\bibitem{ref10} Agent-C. arXiv:2512.23738.
\bibitem{ref11} Pro2Guard. arXiv:2508.00500.
\bibitem{ref12} ProbGuard. arXiv:2602.19844.
\bibitem{ref13} LlamaFirewall. arXiv:2505.03574.
\bibitem{ref14} MAST: multi-agent system failure taxonomy. arXiv:2503.13657.
\bibitem{ref15} AgentOps. arXiv:2411.05285.
\bibitem{ref16} Wasted-computation diagnosis in multi-agent systems.
\bibitem{ref17} LumiMAS. AAMAS 2026. arXiv:2508.12412.
\bibitem{ref18} DiLLS. CHI 2026. arXiv:2602.05446.
\bibitem{ref19} Graph Harness. arXiv:2604.11378.
\bibitem{ref20} SHIELDA. arXiv:2508.07935.
\bibitem{ref21} Learning to Interrupt. arXiv:2604.06452.
\bibitem{ref22} LangGraph interrupt documentation.
\bibitem{ref23} Kephart, Chess. The vision of autonomic computing. IEEE Computer, 2003.
\bibitem{ref24} Weyns. Introduction to self-adaptive systems. 2020.
\bibitem{ref25} Nascimento et al. arXiv:2307.06187.
\bibitem{ref26} ACM TAAS roadmap. DOI 10.1145/3686803.
\bibitem{ref27} Anthropic Engineering. How we built our multi-agent research system.
\bibitem{ref28} Mullick, T\"uz\"un. [Prior work by the authors; citation to be completed at submission.]
\bibitem{ref29} Mullick, T\"uz\"un. [Prior work by the authors; citation to be completed at submission.]
\bibitem{ref30} Zep temporal knowledge graphs. arXiv:2501.13956.
\bibitem{ref31} GAIA. arXiv:2311.12983.
\bibitem{ref32} tau-bench. arXiv:2406.12045.
\bibitem{ref34} Ferrando, Cardoso. RVPLAN: a general purpose framework for replanning using runtime verification. VORTEX@ISSTA 2021, pp. 22--25. DOI 10.1145/3464974.3468447.
\bibitem{ref35} Bozzano, Cimatti, Roveri, Tchaltsev. A comprehensive approach to on-board autonomy verification and validation. IJCAI 2011, pp. 2398--2403.
\bibitem{ref36} Bensalem, Havelund, Orlandini. Verification and validation meet planning and scheduling. STTT 16(1):1--12, 2014.
\bibitem{ref37} Ferrando, Cardoso. Runtime monitoring of action specifications for replanning in classical planning. VORTEX 2025.
\bibitem{ref38} OpenAI. GPT-5.5 model reference and API pricing. developers.openai.com/api/docs/models/gpt-5.5 and /api/docs/pricing (model \texttt{gpt-5.5-2026-04-23}; \$5.00/\$0.50/\$30.00 per 1M input/cached/output tokens; verified 2026-07-01).
\end{thebibliography}

\end{document}
